% Definition des grundsätzlichen Aufbaus
\documentclass[11pt,a4paper]{article}
\usepackage[utf8]{inputenc}
\usepackage[german]{babel}
\usepackage[a4paper,top=30mm,bottom=20mm,left=40mm,right=25mm]{geometry}
\usepackage[onehalfspacing]{setspace}
\usepackage{lmodern}

% Symbole für mathematische Formeln
\usepackage{amsmath}
\usepackage{amsfonts}
\usepackage{amssymb}

% Tabellen und Grafiken
\usepackage{booktabs}
\usepackage{float}
\usepackage{graphicx}
\usepackage{hyperref}

% Sonstige Anpassungen
\usepackage[printonlyused]{acronym}
\setlength\belowcaptionskip{10pt}

% Literaturverzeichnis
\usepackage{csquotes}
\usepackage[style=apa, backend=biber]{biblatex}
\addbibresource{references.bib}

% Beginn des Dokuments
\begin{document}

% Titelseite
\setcounter{page}{0} % nummeriert Titelseite in PDF mit Seite 0
	\begin{titlepage}
		\begin{center}
			{\LARGE\scshape Universität Hamburg}\\*[5mm]
      		{\normalsize Fakultät für Betriebswirtschaft\\
			Institut für Logistik, Transport und Produktion}\\*[3cm]
      		{\Large\scshape Bachelor/Masterarbeit/Seminararbeit}\\*[12mm]
      		{\bfseries\Large\scshape Titel der Bachelor/Masterarbeit}\\*[12mm] % hier Titel einfügen #############
      		{\rmfamily Freie wissenschaftliche Arbeit\\
           	zur Erlangung des Grades "`Bachelor/Master of Science"'\\
           	im Studiengang \ldots} \\*[12mm]
      		{\normalsize\rmfamily Hamburg, \ldots} % hier Datum einfügen #############
		\end{center}
		\vspace{8cm}
		\begin{tabbing}
			Fachsemester: \= \hspace{70mm}\= \kill
 			eingereicht von: \>\> betreut von: \\*[2mm]
      		\textbf{NAME, VORNAME} \>\>  Prof. Dr. Knut Haase\\*[2mm] % hier Name einfügen #############
      		Matrikel-Nr.: \> XXXXXXX \>  \\*[2mm] % hier Matrikelnummer und Betreuer einfügen #############
      		Fachsemester: \> X \\*[2mm] % hier Fachsemester einfügen #############
      		E-Mail: vorname.name@studium.uni-hamburg.de % hier E-Mail-Adresse einfügen #############
		\end{tabbing}
  \end{titlepage}

% Inhaltsverzeichnis
\newpage
\pagenumbering{Roman}
\tableofcontents

% Abbildungsverzeichnis
\newpage
\addcontentsline{toc}{section}{Abbildungsverzeichnis}
\listoffigures

% Tabellenverzeichnis
\newpage
\addcontentsline{toc}{section}{Tabellenverzeichnis}
\listoftables

% Abkürzungsverzeichnis
% Einträge hinzufügen oder entfernen. Alphabetisch sortieren.
% Bei erster Verwendung im Text ausschreiben.
\newpage
\addcontentsline{toc}{section}{Abkürzungsverzeichnis}
\section*{Abkürzungsverzeichnis}
\begin{acronym}[VRP]
	\acro{LLM}{Large Language Model}
	\acro{OR}{Operations Research}
	\acro{VRP}{Vehicle Routing Problem (Tourenplanungsproblem)}
\end{acronym}

% Symbolverzeichnis
% Einträge hinzufügen oder entfernen. Alphabetisch sortieren.
\newpage
\addcontentsline{toc}{section}{Symbolverzeichnis}
\section*{Symbolverzeichnis}
\begin{tabular}{p{0.08\textwidth} p{0.92\textwidth}}
	$c_{ij}$ & Transportkosten von Standort $i$ nach Standort $j$\\
	$d_i$ & Bedarf des Kunden $i \in \mathcal{I}$\\
	$\mathcal{I}$ & Menge der Kunden\\
	$\mathcal{J}$ & Menge aller Standorte (Kunden und Depot)\\
	$x_{ij}$ & Binärvariable, 1 falls ein Fahrzeug von $i$ nach $j$ fährt
\end{tabular}

\newpage
\pagenumbering{arabic}

% ============================================================
%   VERWENDUNG DIESER VORLAGE
%
%   1. Titelseite oben anpassen (############# Markierungen beachten)
%   2. Literaturverweise in references.bib eintragen (Zotero + Better BibTeX)
%   3. Inhalte unten in LaTeX schreiben
%   4. Kompilieren: pdflatex → biber → pdflatex → pdflatex
%
%   KURZREFERENZ
%   Zitate:        \parencite{key} → (Autor, Jahr)
%                  \textcite{key}  → Autor (Jahr)
%                  \parencite{a,b} → mehrere
%                  \parencite[S.~42]{key} → mit Seite
%   Querverweise:  \ref{fig:label}  \ref{tab:label}  \ref{eq:label}  \ref{sec:label}
%   Abbildungen:   \begin{figure}[H] ... \caption{...} \label{fig:...} \end{figure}
%   Tabellen:      \begin{table}[H] ... \caption{...} \label{tab:...} \end{table}
%   Gleichungen:   \begin{align} ... \label{eq:...} \end{align}
%   Abschnitte:    \section{...} \label{sec:...}
%
%   FORMATIERUNG — GRUNDLAGEN
%   Fett:          \textbf{fetter Text}
%   Kursiv:        \textit{kursiver Text}   oder   \emph{hervorgehoben}
%   Unterstrichen: \underline{unterstrichener Text}
%   Schreibmasch.: \texttt{Code oder Dateiname}
%   Aufzählung:    \begin{itemize}
%                    \item Punkt 1
%                    \item Punkt 2
%                    \begin{itemize}
%                      \item Unterpunkt (verschachtelt)
%                    \end{itemize}
%                  \end{itemize}
%   Nummerierung:  \begin{enumerate}
%                    \item Erstens
%                    \item Zweitens
%                  \end{enumerate}
%   Fußnoten:      Text\footnote{Das ist eine Fußnote.}
%   Links:         \href{https://example.com}{Linktext}  (hyperref-Paket)
%   Abstände:      \vspace{1cm}  (vertikal)
%                  \hspace{1cm}  (horizontal)
%                  ~             (geschütztes Leerzeichen, z.B. Abbildung~\ref{})
%                  \,            (schmales Leerzeichen, z.B. 16\,GB)
%   Gedankenstrich: -- (für Bereiche, z.B. Seiten 5--10)
%
%   HÄUFIGE FEHLER
%   - Jede Abbildung und Tabelle muss im Text referenziert werden,
%     *bevor* sie erscheint.
%   - Abbildungen haben Beschriftungen darunter (\caption nach \includegraphics),
%     Tabellen darüber (\caption vor \begin{tabular}).
%   - Sätze nicht mit einem Symbol beginnen ("$X$ ist..." →
%     "Die Variable $X$ ist...").
%   - Notation konsistent halten — nicht zwischen $X_{ij}$ und $X(i,j)$
%     wechseln.
%   - Keine verwaisten Abschnitte (ein einzelner \subsection in einem
%     \section).
%   - Jeder Abschnitt/Unterabschnitt sollte mindestens eine Seite lang sein.
%     Wenn nicht, mit einem anderen Abschnitt zusammenfassen.
%   - Abkürzungen bei der ersten Verwendung mit \ac{} ausschreiben,
%     danach konsequent die Kurzform verwenden.
%   - Wikipedia und Vorlesungsfolien nicht zitieren — die Originalquelle
%     suchen.
%   - Gedankenstriche (--) für Bereiche verwenden (Seiten 5--10), keine
%     Bindestriche (-).
%   - Zahlen eins bis zwölf im Fließtext ausschreiben, ab 13 Ziffern
%     verwenden.
%   - Schmales Leerzeichen vor Einheiten: 16\,GB, 245\,s.
%   - Geschütztes Leerzeichen vor \ref verwenden: Abbildung~\ref{fig:...}
% ============================================================

% SCHREIBHILFE: Einführung
% Strukturieren Sie die Einleitung als Trichter:
%   1. Breiter Kontext — warum ist dieses Feld relevant?
%   2. Konkretes Problem — was ist die spezifische Herausforderung?
%   3. Forschungslücke — was fehlt in der bisherigen Literatur?
%   4. Eigener Beitrag — was leistet diese Arbeit?
%   5. Aufbau der Arbeit — kurz die Struktur beschreiben.
% Tipp: Schreiben Sie diesen Abschnitt zuletzt, wenn Sie Ihren tatsächlichen
% Beitrag kennen. Er sollte beantworten: Warum sollte der Leser sich
% dafür interessieren? Was ist neu?
\section{Einführung}
\label{sec:einfuehrung}
Die Tourenplanung macht einen erheblichen Anteil der Distributionskosten im Straßengüterverkehr aus. Seit das \ac{VRP} erstmals von \textcite{dantzig1959truck} formalisiert wurde, haben es Fortschritte bei den Lösungsverfahren ermöglicht, immer umfangreichere Probleminstanzen zu lösen \parencite{desaulniers2005column}. Die vorliegende Arbeit befasst sich mit dem Problem \ldots

Abbildung \ref{fig:beispiel} zeigt eine schematische Übersicht der Problemstruktur.

\begin{figure}[H]
	\centering
	\includegraphics[width=0.8\textwidth]{images/Image.pdf}
	\caption{Ersetzen Sie dies durch Ihre eigene Abbildung.}
	\label{fig:beispiel}
\end{figure}

Die Arbeit ist wie folgt gegliedert. Abschnitt \ref{sec:literatur} gibt einen Überblick über die relevante Literatur. Abschnitt \ref{sec:modell} stellt das mathematische Modell vor, wobei Abschnitt \ref{sec:annahmen} die zugrundeliegenden Annahmen und Abschnitt \ref{sec:formulierung} die formale Formulierung enthält. Abschnitt \ref{sec:rechenstudie} präsentiert die Rechenstudie, Abschnitt \ref{sec:diskussion} diskutiert die Befunde, und Abschnitt \ref{sec:zusammenfassung} fasst die Ergebnisse zusammen und gibt einen Ausblick.

% SCHREIBHILFE: Literaturüberblick
% - Quellen thematisch oder methodisch gruppieren, nicht chronologisch.
% - Nicht nur zusammenfassen — Ansätze vergleichen, Stärken und Schwächen
%   aufzeigen und die Lücke identifizieren, die Ihre Arbeit füllt.
% - Jede zitierte Arbeit sollte einen Bezug zu Ihrer Forschungsfrage haben.
% - Am Ende überleiten zum eigenen Beitrag: Was bleibt ungelöst?
% - \parencite{key} für Klammerzitate und \textcite{key} für narrative Zitate.
\section{Literaturüberblick}
\label{sec:literatur}
Das \ac{VRP} wird seit der grundlegenden Arbeit von \textcite{dantzig1959truck} umfassend untersucht. Im Laufe der Jahrzehnte wurden zahlreiche exakte und heuristische Ansätze vorgeschlagen \parencite{toth2014vehicle}.

Spaltengenerierung hat sich als besonders effektive Technik zur Lösung umfangreicher Touren- und Ablaufplanungsprobleme erwiesen \parencite{desaulniers2005column}. Zudem werden seit einiger Zeit Methoden des maschinellen Lernens eingesetzt, um die Leistung von Heuristiken zu verbessern \parencite{bengio2021machine}.

Ein umfassender Überblick über Lösungsmethoden geht über den Rahmen dieser Arbeit hinaus. Für eine detaillierte Darstellung sei auf \textcite{toth2014vehicle} verwiesen.

% SCHREIBHILFE: Modell
% Kurzen Überblick über den Modellierungsansatz geben, bevor Sie ins Detail
% gehen. Aufteilen in Annahmen (was vereinfacht wird) und Formulierung
% (die Mathematik).
\section{Modell}
\label{sec:modell}
Dieser Abschnitt stellt das mathematische Modell für das oben beschriebene Problem vor. Zunächst werden die wesentlichen Annahmen erläutert, anschließend folgt die formale Formulierung.

% SCHREIBHILFE: Annahmen
% - Jede Annahme muss begründet werden: Warum ist sie vertretbar?
% - Einschränkungen benennen: Was würde sich ändern, wenn eine Annahme
%   gelockert wird?
% - Von der grundlegendsten zur spezifischsten Annahme ordnen.
\subsection{Annahmen}
\label{sec:annahmen}
Die Modellierung des Problems erfordert einige vereinfachende Annahmen:
\begin{itemize}
	\item Ein einzelnes Depot dient als Start- und Endpunkt für alle Fahrzeuge.
	\item Jeder Kunde $i \in \mathcal{I}$ hat einen bekannten Bedarf $d_i$, der vollständig erfüllt werden muss.
	\item Die Flotte besteht aus identischen Fahrzeugen mit einheitlicher Kapazität.
	\item Die Transportkosten $c_{ij}$ zwischen den Standorten $i$ und $j$ sind symmetrisch und erfüllen die Dreiecksungleichung.
\end{itemize}

% SCHREIBHILFE: Formulierung
% - Alle Notation definieren, bevor sie verwendet wird.
% - Konsistente Konventionen: Mengen = kalligraphisch ($\mathcal{J}$),
%   Variablen = Großbuchstaben ($X_j$), Parameter = Kleinbuchstaben ($b_i$).
% - Jede Nebenbedingung muss im Text erläutert werden — nicht nur
%   Gleichungen auflisten.
% - Nebenbedingungen per Label referenzieren:
%   "Nebenbedingung (\ref{eq:zuweisung}) stellt sicher..."
%
% TIPPS ZUR TABELLENGESTALTUNG
% - Booktabs-Stil verwenden (\toprule, \midrule, \bottomrule) — keine
%   vertikalen Linien.
% - Zahlenspalten rechtsbündig (r), Textspalten linksbündig (l) ausrichten.
% - Einheiten in Spaltenüberschriften, nicht in jeder Zelle.
% - Tabellen kompakt halten — unnötige Spalten vermeiden.
\subsection{Formulierung}
\label{sec:formulierung}
Mengen werden kalligraphisch gesetzt ($\mathcal{J}$), Variablen als Großbuchstaben ($X_j$, $j \in \mathcal{J}$) und Parameter sowie Indizes als Kleinbuchstaben ($b_i$).
% HINWEIS: Eine vollständige VRP-Formulierung würde zusätzlich Kapazitäts-
% und Subtour-Eliminationsbedingungen enthalten. Sie werden hier weggelassen,
% um das Beispielmodell kurz zu halten.
Das Modell ist wie folgt definiert:
\begin{align}
	\label{eq:zielfunktion}
	\min \sum_{i \in \mathcal{J}} \sum_{j \in \mathcal{J}} c_{ij} \cdot x_{ij} & \\
	\intertext{unter den Nebenbedingungen}
	\label{eq:zuweisung}
	\sum_{j \in \mathcal{J}} x_{ij} = 1 && \forall i \in \mathcal{I} \\
	\label{eq:fluss}
	\sum_{i \in \mathcal{J}} x_{ij} = 1 && \forall j \in \mathcal{I} \\
	\label{eq:binaer}
	x_{ij} \in \{0, 1\} && \forall i, j \in \mathcal{J}
\end{align}
Die Zielfunktion (\ref{eq:zielfunktion}) minimiert die gesamten Transportkosten. Nebenbedingung (\ref{eq:zuweisung}) stellt sicher, dass genau ein Fahrzeug von jedem Kunden abfährt, während (\ref{eq:fluss}) sicherstellt, dass genau ein Fahrzeug bei jedem Kunden ankommt. (\ref{eq:binaer}) definiert die binären Entscheidungsvariablen.

% SCHREIBHILFE: Rechenstudie
% - Hardware (CPU, RAM), Software (Solver, Version) und Sprache angeben.
% - Instanzquellen beschreiben oder erklären, wie Testdaten erzeugt wurden.
% - Immer gegen eine Baseline, Benchmark oder exakte Lösung vergleichen.
% - Sowohl Lösungsqualität (Gap, Zielfunktionswert) als auch Rechenzeit
%   berichten.
% - Trends über Instanzgrößen hinweg diskutieren, nicht nur einzelne Zahlen
%   auflisten.
% - Zeitlimits, Anzahl der Durchläufe oder andere experimentelle Parameter
%   erwähnen.
\section{Rechenstudie}
\label{sec:rechenstudie}
Zur Bewertung des vorgeschlagenen Ansatzes führen wir eine Reihe von Rechenexperimenten durch. Alle Experimente wurden auf einem Rechner mit Intel Core i7 Prozessor und 16\,GB RAM durchgeführt. Das Modell wurde in Python unter Verwendung von Gurobi 11.0 als Solver implementiert.

Tabelle \ref{tab:ergebnisse} zeigt die Rechenergebnisse im Vergleich von exakter und heuristischer Lösung.
\begin{table}[H]
	\centering
	\caption{Rechenergebnisse im Vergleich von exakter und heuristischer Lösung}
	\label{tab:ergebnisse}
	\begin{tabular}{lcrrrr}
		\toprule
		Instanz & $|\mathcal{I}|$ & Optimal & Heuristik & Gap (\%) & Zeit (s)\\
		\midrule
		Klein-1  & 25  &   245,3 &   248,1 & 1,14 &    0,8\\
		Klein-2  & 25  &   312,7 &   318,4 & 1,82 &    1,1\\
		Mittel-1 & 50  & 1.340,5 & 1.382,9 & 3,16 &   12,4\\
		Mittel-2 & 50  & 1.128,3 & 1.167,2 & 3,45 &   15,7\\
		Groß-1   & 100 & 3.450,1 & 3.621,3 & 4,96 &  238,7\\
		Groß-2   & 100 & 4.217,8 & 4.485,6 & 6,35 &  310,2\\
		\bottomrule
	\end{tabular}
\end{table}

Die Ergebnisse in Tabelle \ref{tab:ergebnisse} zeigen, dass der heuristische Ansatz nahezu optimale Lösungen innerhalb vertretbarer Rechenzeiten findet. Für kleine Instanzen bleibt die Optimalitätslücke unter 2\%, während sie für größere Instanzen auf etwa 5--6\% ansteigt.

% SCHREIBHILFE: Diskussion
% - Ergebnisse im Kontext der Forschungsfrage interpretieren.
% - Ergebnisse mit bestehender Literatur vergleichen: Sind sie konsistent?
%   Wo unterscheiden sie sich, und warum?
% - Limitationen ehrlich diskutieren — jede Studie hat welche.
% - Beobachtung (die Daten zeigen X) von Interpretation (dies deutet auf Y
%   hin) unterscheiden.
% - Unerwartete Ergebnisse ansprechen, nicht ignorieren.
% - Keine völlig neuen Themen einführen. Neue Forschungsrichtungen
%   gehören in den Ausblick der Zusammenfassung.
\section{Diskussion}
\label{sec:diskussion}
Die Ergebnisse deuten darauf hin, dass die vorgeschlagene Heuristik für kleine und mittlere Instanzen gut funktioniert, jedoch bei größeren Problemgrößen zunehmende Lücken aufweist. Dies ist konsistent mit den Beobachtungen von \textcite{toth2014vehicle}, die ein ähnliches Skalierungsverhalten für konstruktive Heuristiken bei kapazitätsbeschränkten Tourenplanungsproblemen feststellten.

Eine wesentliche Einschränkung dieser Studie ist die Verwendung symmetrischer Kostenmatrizen, die in realen Logistiknetzwerken mit Einbahnstraßen oder zeitabhängigen Fahrzeiten möglicherweise nicht zutreffen. Zudem wurden die Testinstanzen zufällig generiert und nicht aus realen Daten abgeleitet, was die Übertragbarkeit der Ergebnisse einschränkt.

% SCHREIBHILFE: Zusammenfassung
% - Zusammenfassung (was wurde gemacht und gefunden) von Ausblick
%   (zukünftige Forschung) trennen.
% - Keine neuen Ergebnisse oder Analysen hier einführen.
% - Zukünftige Forschungsrichtungen sollten konkret und umsetzbar sein,
%   nicht vage ("weitere Forschung ist nötig").
% - Nicht übertreiben: "deutet darauf hin" statt "beweist",
%   "legt nahe" statt "zeigt" — es sei denn, ein formaler Beweis liegt vor.
\section{Zusammenfassung und Ausblick}
\label{sec:zusammenfassung}
Die vorliegende Arbeit hat ein mathematisches Modell und einen heuristischen Lösungsansatz für das Tourenplanungsproblem vorgestellt. Die Rechenstudie legt nahe, dass die vorgeschlagene Methode nahezu optimale Lösungen in einem Bruchteil der Rechenzeit des exakten Verfahrens findet.

Aus dieser Arbeit ergeben sich zwei Richtungen für zukünftige Forschung. Das Modell könnte um Zeitfenster und heterogene Flotten erweitert werden. Darüber hinaus könnte die Heuristik durch Methoden des maschinellen Lernens verbessert werden, um die Qualität der Startlösungen zu erhöhen.

% Literaturverzeichnis ausgeben
\newpage
\pagenumbering{Roman}
\setcounter{page}{6} % ggf. an die tatsächliche Seitenzahl anpassen
\addcontentsline{toc}{section}{Literatur}
\printbibliography

\newpage
\addcontentsline{toc}{section}{Eidesstattliche Erklärung}
\section*{Eidesstattliche Erklärung}
Ich versichere an Eides statt, dass ich die vorstehende Arbeit selbstständig und ohne fremde Hilfe angefertigt und mich anderer als der im beigefügten Verzeichnis angegebenen Hilfsmittel nicht bedient habe. Alle Stellen, die wörtlich oder sinngemäß aus Veröffentlichungen übernommen wurden, sind als solche kenntlich gemacht. Alle Internetquellen sind der Arbeit beigefügt. Des Weiteren versichere ich, dass ich die Arbeit vorher nicht in einem anderen Prüfungsverfahren eingereicht habe und dass die eingereichte schriftliche Fassung der auf dem elektronischen Speichermedium entspricht.

\vspace{2cm}
\begin{flushleft}
Hamburg, Datum \\
\vspace{0.5cm}
Vorname Nachname \end{flushleft}

\newpage
\section*{Verwendung von KI-Werkzeugen}
\addcontentsline{toc}{section}{Verwendung von KI-Werkzeugen}

% Geben Sie an, welche KI-Werkzeuge Sie verwendet haben und wie.
% Seien Sie konkret: welches Tool, für welche Aufgaben, und wie Sie die Ergebnisse geprüft haben.
% Beispiele:
%   - "ChatGPT: Verwendet zur Strukturierung des Literaturüberblicks"
%   - "Claude: Verwendet zum Debugging von Python-Code in der Rechenstudie"
%   - "GitHub Copilot: Verwendet zur Code-Vervollständigung in der Implementierung"
%   - "Grammarly: Verwendet zum Korrekturlesen und zur Verbesserung von Abschnitt 3"
% Alle KI-generierten Inhalte müssen eigenständig geprüft, verifiziert und überarbeitet werden.

Die folgenden KI-Werkzeuge wurden bei der Erstellung dieser Arbeit eingesetzt:

\begin{itemize}
	\item \textbf{Tool 1}: Verwendet für \ldots
	\item \textbf{Tool 2}: Verwendet für \ldots
	\item \textbf{Tool 3}: Verwendet für \ldots
\end{itemize}

\vspace{2cm}
\begin{flushleft}
Hamburg, Datum \\
\vspace{0.5cm}
Vorname Nachname \end{flushleft}
\end{document}
